\chapter{INTRODUCTION}

\section{1.1 BACKGROUND}
The contemporary paradigm of municipal waste disposal and material recovery requires a fundamental overhaul to meet the demands of rapid global urbanization. Historically, the scrap metal and recyclable materials sector has operated as an unstructured, fragmented, and largely informal economy. Individuals wishing to dispose of recyclable commodities are often forced to negotiate with local, unregulated collectors, leading to a disparity in material valuation and widespread distrust. In an era dominated by digital transformation, the lack of an organized, data-driven infrastructure for waste recovery is a prominent bottleneck in achieving sustainable smart city objectives. The integration of Artificial Intelligence (AI) and cloud computing offers an unprecedented opportunity to formalize this sector. ScrapKart AI is envisioned to address this systemic issue by providing a digital marketplace that connects consumers with verified scrap collectors.

\section{1.2 MOTIVATION}
The primary catalyst for developing ScrapKart AI stems from the urgent ecological necessity to divert recyclable materials away from overburdened landfills. Despite high public willingness to recycle, the logistical friction and lack of standardized financial incentives severely hamper participation rates. Furthermore, independent scrap collectors often suffer from inefficient routing, wasting fuel and time searching for materials. Observing the success of digital marketplaces in ride-sharing and food delivery, it became evident that applying similar logistical frameworks, augmented by predictive machine learning algorithms for price estimation, could revolutionize the recycling ecosystem. 

\section{1.3 PROBLEM DEFINITION}
Traditional scrap collection ecosystems suffer from severe limitations, such as inconsistent pricing, inefficient scheduling, and lack of digital integration. These limitations reduce user participation and negatively impact recycling efficiency. The proposed system, ScrapKart AI, addresses these challenges by enabling:
\begin{itemize}
    \item Real-time AI-based scrap classification from uploaded images.
    \item Dynamic price calculation based on real-time market trends.
    \item Location-based dispatching of nearest scrap collectors.
    \item Transparent digital ledger records for all recycling transactions.
\end{itemize}
The primary problem is to design a system capable of accurately classifying various types of recyclable scrap from raw image data while simultaneously optimizing the logistics network for everyday consumers.

\section{1.4 ORGANIZATION OF REPORT}
The remainder of this report is organized into structured chapters that explore all components of ScrapKart AI in detail.\\
\textbf{Chapter 1} - presents the introduction, background, and problem definition of the system.\\
\textbf{Chapter 2} - reviews existing research on smart waste management and computer vision algorithms.\\
\textbf{Chapter 3} - provides the Software Requirements Specification (SRS), including functional, non-functional, and system requirements.\\
\textbf{Chapter 4} - describes the system architecture, data flow diagrams, ER diagrams, and UML models.\\
\textbf{Chapter 5} - outlines the technical specifications and technologies used in the project.\\
\textbf{Chapter 6} - details the COCOMO project planning, cost estimation, and development schedule.\\
\textbf{Chapter 7} - discusses software implementation, database design, and core algorithms.\\
\textbf{Chapter 8} - highlights the software testing procedures and comprehensive test cases.\\
\textbf{Chapter 9} - showcases experimental results, system evaluation, and observations.\\
\textbf{Chapter 10} - concludes the report and highlights future expansion possibilities for smart city recycling technologies.
